\unit{Hydrogenic atoms and radial wavefunctions}
\unitref{Week 8; 2023 Q2; Assignment 3 Q1; Assignment 4 Q3}
\subsection{The one-electron Schr\"odinger equation}
For an electron of reduced mass \(\mu\)
around a point nucleus of charge \(+Ze\),
\[
\left[-\frac{\hbar^2}{2\mu}\nabla^2
-\frac{Ze^2}{4\pi\epsilon_0r}\right]
\psi(\mathbf r)=E\psi(\mathbf r).
\]
The Coulomb potential depends only on
the radius \(r\), so angular momentum is
conserved. Separate variables into
\[
\boxed{\psi_{nlm}(\mathbf r)=R_{nl}(r)
Y_{lm}(\theta,\phi).}
\]
The spherical harmonics are eigenfunctions
of \(L^2,L_z\) with
\[
L^2Y_{lm}=\hbar^2l(l+1)Y_{lm},\qquad
L_zY_{lm}=\hbar mY_{lm}.
\]
The labels have allowed values
\(n=1,2,\ldots\), \(l=0,\ldots,n-1\),
\(m=-l,\ldots,l\).

\subsection{The centrifugal barrier}
Define \(u(r)=rR_{nl}(r)\).
Then the radial Schr\"odinger equation becomes
\[
\boxed{-\frac{\hbar^2}{2\mu}\frac{d^2u}{dr^2}
+\left[V(r)+\frac{\hbar^2l(l+1)}{2\mu r^2}\right]u
=Eu.}
\]
The term proportional to \(1/r^2\)
is the effective centrifugal barrier.
For \(l>0\), the radial wavefunction
is suppressed near the nucleus;
\(s\)-states with \(l=0\) penetrate it
much more strongly.
The normalisation can be written
\[
\int_0^\infty|R_{nl}|^2r^2dr
=\int_0^\infty|u_{nl}|^2dr=1.
\]
Switching from \(R\) to \(u\) changes the
\emph{measure} in every radial matrix element.

\subsection{Why hydrogen energies depend only on \(n\)}
The nonrelativistic Coulomb levels are
\[
\boxed{E_n=-\frac{\mu Z^2e^4}
{2(4\pi\epsilon_0)^2\hbar^2n^2}
\approx-13.6\,{\rm eV}\,
\frac{\mu}{m_e}\frac{Z^2}{n^2}.}
\]
The \(n\)-only degeneracy is a special
symmetry of the pure Coulomb potential.
Screening in a many-electron atom,
relativity, magnetic fields and finite
nuclear size lift different parts
of this degeneracy.

\begin{worked}{Most probable radius is not mean radius}
For hydrogen \(1s\),
\(\psi_{1s}=e^{-r/a_0}/\sqrt{\pi a_0^3}\).
The probability in a spherical shell
between \(r\) and \(r+dr\) is
\[
P(r)dr=4\pi r^2|\psi|^2dr
=\frac{4r^2}{a_0^3}e^{-2r/a_0}dr.
\]
Differentiate \(\ln P(r)\):
\(2/r-2/a_0=0\), so
\(\boxed{r_{\rm most\ probable}=a_0}\).
But
\[
\langle r\rangle=\frac4{a_0^3}
\int_0^\infty r^3e^{-2r/a_0}dr
=\boxed{3a_0/2}.
\]
The different answers are not contradictory:
one describes the peak of a distribution;
the other its average.
\end{worked}

\subsection{Finite nuclear size is a perturbation}
A nucleus has a finite radius and a spread-out
charge distribution. The point Coulomb potential
\(-Z/r\) is too singular near the origin.
For a uniform nuclear sphere of radius \(R\)
in atomic units,
\[
V_{\rm finite}(r)=
-\frac{Z}{2R}\left(3-\frac{r^2}{R^2}\right)
\quad(r<R).
\]
The perturbation is
\[
\Delta V=V_{\rm finite}-V_{\rm point}
=Z\left(\frac1r-\frac{3}{2R}
+\frac{r^2}{2R^3}\right).
\]
First-order perturbation theory gives
\[
\boxed{\Delta E_{nl}=\int_0^R
|R_{nl}(r)|^2r^2\Delta V(r)\,dr.}
\]
If \(R\) is small and the \(s\)-state
wavefunction is approximately constant
inside the nucleus, one obtains
\[
\boxed{\Delta E_{ns}\simeq
\frac{2Z^4R^2}{5n^3}\quad\text{(Hartree)}.}
\]
The positive shift makes the state less
strongly bound. It is particularly
large for orbitals with nonzero density
at the origin.

\subsection{How numerical radial solutions work}
For a screened or model potential \(V(r)\),
the radial equation is often discretised.
At grid spacing \(h\),
\[
u''(r_i)\approx
\frac{u(r_{i-1})-2u(r_i)+u(r_{i+1})}{h^2}.
\]
The Hamiltonian becomes a sparse
tridiagonal matrix. Diagonalising it
yields energies and radial eigenvectors.
Check normalisation, radial nodes and
stability as \(h\) decreases and
the outer boundary moves outward.
The sodium lifetime assignment uses
exactly this kind of numerical workflow.

\begin{checkpoint}
Hydrogenic \(n,l,m\) label wavefunctions,
not miniature classical orbits.
For a radial integral, always decide
whether you are using \(R(r)\) or
\(u(r)=rR(r)\) before writing the measure.
\end{checkpoint}
